% Generated draft: requires mathematical and source review.
\documentclass[11pt]{article}
\usepackage[margin=25mm]{geometry}
\usepackage{fontspec}
\setmainfont{Latin Modern Roman}
\usepackage{amsmath,amssymb,mathtools,unicode-math}
\setmathfont{Latin Modern Math}
\usepackage{enumitem,needspace,xurl,hyperref}
\setlength{\parindent}{0pt}
\setlength{\parskip}{5pt}
\setlength{\emergencystretch}{3em}
\newcommand{\N}{\mathbb{N}}
\newcommand{\Z}{\mathbb{Z}}
\newcommand{\Q}{\mathbb{Q}}
\newcommand{\R}{\mathbb{R}}
\newcommand{\C}{\mathbb{C}}
\newcommand{\F}{\mathbb{F}}
\newcommand{\A}{\mathbb{A}}
\newcommand{\PP}{\mathbb{P}}
\newcommand{\E}{\mathbb{E}}
\newcommand{\eps}{\varepsilon}
\newcommand{\dd}{\,\mathrm{d}}
\begin{document}
\section*{UnsolvedMath ID: 801177}
\textbf{Source problem code:} catalogue entry
\subsection{Definitions and mathematical statement}
Let \(e\geq 1\) be an integer. Let \(\Delta\) be a finite abstract simplicial complex of dimension \(2e-1\), so its largest faces have \(2e\) vertices. A minimal non-face is a set of vertices that is not a face although every proper subset is a face. The complex is flag if every minimal non-face has exactly two vertices.

To make the coefficient convention explicit, fix a field \(k\). For a face \(\sigma\in\Delta\), define
\[
\operatorname{lk}_{\Delta}(\sigma)=\{\tau\in\Delta:\tau\cap\sigma=\varnothing,\ \tau\cup\sigma\in\Delta\}.
\]
Use the standard homological formulation of Gorenstein* over \(k\): the complex is pure of dimension \(2e-1\), and for every face \(\sigma\), including the empty face,
\[
\widetilde H_i(\operatorname{lk}_{\Delta}(\sigma);k)\cong
\begin{cases}
k,&i=2e-1-|\sigma|,\\
0,&i\neq 2e-1-|\sigma|.
\end{cases}
\]
Here reduced simplicial homology includes degree \(-1\), with \(\widetilde H_{-1}(\{\varnothing\};k)\cong k\).

Let \(f_{j-1}\) count the faces with \(j\) vertices, with \(f_{-1}=1\). Define the \(h\)-vector by
\[
h_{\Delta}(t)=\sum_{i=0}^{2e}h_i t^i
=\sum_{j=0}^{2e}f_{j-1}t^j(1-t)^{2e-j}.
\]
The proposed assertion is that every such Gorenstein* flag complex satisfies
\[
(-1)^e h_{\Delta}(-1)
=(-1)^e\sum_{i=0}^{2e}(-1)^i h_i\geq 0.
\]
The supplied record does not specify a coefficient field or its quantification. The formulation above is relative to a fixed field \(k\); the intended field scope requires confirmation. Finiteness is the conventional assumption underlying the supplied face-counting \(h\)-vector. The record also leaves the range of \(e\) implicit; the statement above covers positive odd dimensions. If \(e=0\) and the dimension \(-1\) complex \(\{\varnothing\}\) is admitted, its \(h\)-polynomial is \(1\), and the displayed inequality holds.
\subsection{Short English statement}
The Charney--Davis conjecture in Stanley's Gorenstein* flag formulation asks whether the alternating sum of the \(h\)-vector has sign \((-1)^e\) or is zero in dimension \(2e-1\). The coefficient-field scope is unspecified in the supplied record and requires review.
\subsection{Sources}
\begin{itemize}
\item Local UnsolvedMath registry, permanent ID \(801177\), source problem code \texttt{catalogue entry}; supplied complete source record, upstream revision \texttt{d371ee917ac9647fbb06cd6fc529deec6f3dda82}. The registry labels the problem open; this is supplied status metadata, not an independently verified literature assessment.
\item Richard Stanley, \emph{Positivity Problems and Conjectures in Algebraic Combinatorics} (1999), Problem 4. Supplied source URL: \url{https://math.mit.edu/~rstan/papers/problems.pdf}. The supplied provenance records access on 2026-10-16 and extraction marked \texttt{pdf+html\_verified}; these are supplied verification claims and were not independently checked here.
\item The supplied background attributes the conjecture for flag homology spheres to Ruth Charney and Michael Davis, \emph{Pacific Journal of Mathematics} \textbf{171} (1995), 117--137, and its extension to Gorenstein* flag complexes to Stanley. No article title or separate link for the Charney--Davis reference is supplied.
\item The supplied background reports partial results by Leung--Reiner and Reiner--Stanton--Welker, without bibliographic details or descriptions of their scope. These reports are retained as source qualifications and have not been independently verified.
\end{itemize}
\end{document}
